\section{A finite coefficient field for the required theta jets}
\label{jcf-section}

We now pass from the algebraicity in Corollary~\ref{gm-algebraic-jets}
to one number field containing every coefficient used in a fixed
finite column construction. The distinction matters: a collection
of algebraic numbers need not lie in a single number field.
The proof uses the finiteness of positive definite lattice classes,
finite residue quotients, and the fixed bound on the orders of the
required derivatives. This is the finite-generation argument needed
for the submitted coefficient-field construction.\cite{alpo-uniform-source}

\subsection{The data fixed before adjoining coefficients}

Fix the CM field $E$, its CM type $\Sigma$, the finitely many
positive Hermitian spaces, their even ranks, and the normalized
theta family of Sections~\ref{te-section} and \ref{gm-section}.
Fix also the finite list of rational primes $T_u$ used in the
columns, the auxiliary primes $g_i$, and the finite family
$\Xi$ of finite-order character twists. These are arithmetic
data over $E$; they do not depend on an ambient coefficient
field. Let $\mathcal B$ be the finite set of primes of $E$
consisting of the primes above the $T_u,g_i$ and the primes
above the ramification loci of the characters in $\Xi$.
Taking all primes above these rational primes permits every
choice of one prime from a conjugate pair.

Let $m$ be the fixed length of the chain. For $R\subseteq
\mathcal B$ and an integer $b\in\{m,m+1\}$ put
\[
 q_R=\prod_{P\in R}NP,\qquad
                   M_{u,R,b}=T_uq_R^{2b}.                 \tag{JF1}
\]
The exponent $m$ is sufficient for the initial construction;
$m+1$ includes its prescribed increase of local level.
There are finitely many such positive integers. Fix a bound
$d$ on the total order of all derivatives in the finite
list of weights. One can instead take the exact order for
each step; taking their maximum only enlarges a finite set.
Fix the finitely many algebraic direction coefficients
$c_\eta$. The actual construction chooses its eigen-direction
lattices and their generators before the coefficient-field
enlargement; their values are included among the generators
below.

We shall allow all $S\subseteq\mathcal B$ for which the theta
indicator is defined, all positive divisors $c\mid M$ of
one of the integers in (JF1), and all integral self-dual
lattices in the fixed Hermitian spaces. Enlarging the finite
list in this way makes it independent of the later character
realizations and choices of prime representatives.

\subsection{Characters and finite residue quotients}

Let $\mathcal I^-$ be the group of fractional ideals $J$ of
$E$ satisfying $J\bar J=\mathcal O_E$, and let
$\mathcal P^-=\{(u):u\bar u=1\}$. The determinant character
used here satisfies
\[
                  \chi((u))=\prod_{\eta\in\Sigma}\eta(u)^{-1}
                         \quad(u\bar u=1).               \tag{JF2}
\]

\begin{lemma}\label{jcf-character-values}
The group $\mathcal I^-/\mathcal P^-$ is finite. Consequently
all values of $\chi$ belong to one number field.
\end{lemma}
\begin{proof}
The map to the ideal class group has finite image. Its kernel
consists of principal ideals $(z)$ with $z\bar z\in
\mathcal O_{E^+}^{\times,+}$. On that kernel the assignment
\[
 (z)\longmapsto z\bar z\pmod {N_{E/E^+}\mathcal O_E^\times}
\]
is well defined, because replacing $z$ by $z\varepsilon$ changes
its norm by $N\varepsilon$. Its kernel is precisely
$\mathcal P^-$: if $z\bar z=\varepsilon\bar\varepsilon$, then
$z/\varepsilon$ has norm one. The unit quotient on the right
is finite. Indeed the norm subgroup contains the squares of
$\mathcal O_{E^+}^\times$, and the unit theorem makes the
quotient by squares finite. This proves the first assertion.

Choose representatives $J_1,\ldots,J_r$ of this finite group
and let $h$ kill it. For each $i$ there is $u_i\in E^\times$
with $u_i\bar u_i=1$ and $J_i^h=(u_i)$. Formula (JF2) gives
\[
                  \chi(J_i)^h=
                         \prod_\eta\eta(u_i)^{-1}
                          \in\overline{\mathbb Q}^{\times}.
\]
Thus each $\chi(J_i)$ is algebraic. Every ideal $J$ has the
form $J_i(u)$ with $u\bar u=1$. Its value is
$\chi(J_i)\prod_\eta\eta(u)^{-1}$, which lies in the field
generated by the normal closure of $E$ and the finitely many
$\chi(J_i)$. This proves the assertion without adjoining
unboundedly many roots of arbitrary elements of $E$.
\end{proof}

A finite-order twist has finite image, contained in the roots
of unity of its order. Its ramification locus is finite.
Hence all its values, the roots of unity of orders
$(NP)^{2m+2}$ required at the finitely many local levels, and
the finitely many square roots of local norms lie in a number
field. Any prescribed finite list of algebraic frame
coefficients can be adjoined at the same time.
For algebraic Hecke characters with nonzero infinity type,
this statement concerns their normalized finite-idele
values; no assertion is made that their values on arbitrary
archimedean ideles are algebraic.

For two fixed self-dual lattices $L,L'$ write
\[
 H_S(L,L')=\{X\in\operatorname{Hom}_E(V,V'):
                         XL_P\subseteq L'_P\ (P\in S)\}.
\]

\begin{lemma}\label{jcf-finite-indicators}
For positive integers $M,c$ the quotient
\[
                  M^{-1}H_S(L,L')/cH_S(L,L')              \tag{JF3}
\]
is finite. Consequently the theta indicators
$\Theta_{L,L',S,c,x}$ with $x\in M^{-1}H_S(L,L')$ take
only finitely many distinct values as holomorphic functions.
\end{lemma}
\begin{proof}
Put $R_S=\bigcap_{P\in S}\mathcal O_{E,P}\subset E$.
If $S$ is empty, $R_S=E$ and the quotient is zero.
Otherwise $R_S$ is a semilocal Dedekind ring. An invertible
ideal of it is principal: choose, by the Chinese remainder
theorem, an element with its prescribed valuations at the
finitely many maximal ideals. Thus every finitely generated
projective $R_S$-module is free. In particular $H_S(L,L')$
is free of rank $ab$ over $R_S$, where $a,b$ are the ranks
of $L,L'$. Multiplication by $M$ identifies (JF3) with
$H_S/McH_S$. Finally
\[
 R_S/McR_S\simeq
     \prod_{P\in S}\mathcal O_E/P^{v_P(Mc)}
\]
with the factors of exponent zero omitted. The product is
finite. These facts also follow directly by choosing local
bases and applying the Chinese remainder theorem.

The finite theta indicator depends only on the residue class
of $x$ modulo $cH_S$. This is its coset definition: the two
lifts differ by the period of the characteristic function.
Equivalently, their Weyl operators differ by a stabilizer of
the averaged vector, with the same semicharacter phase. The
lift-independence calculation in Section~\ref{csp-section}
proves exactly this equality, including the phase. Thus the
number of distinct functions is bounded by the size of (JF3).
\end{proof}

\subsection{All isometric lattices use the same field}

Let $F_{L,L',S,c,x}$ denote the actual theta indicator divided
by its determinant-character factor. Its mixed derivatives
are the coefficients denoted by \emph{rational indices} in
the source. All of them are algebraic. Indeed the indicator
is a finite linear combination of rational translates of
the normalized theta vector; its coefficients are rational
numbers, square roots of positive integers, and roots of
unity. Corollary~\ref{gm-algebraic-jets}, with the coordinate
directions as in Lemma~\ref{csp-coordinate-units}, gives the
algebraicity of each derivative. Lemma~\ref{jcf-character-values}
permits division by its character factor.

\begin{proposition}\label{jcf-common-field}
There is a number field $K\subset\mathbb C$ containing all
mixed derivatives of order at most $d$ of
$F_{L,L',S,c,x}$ at zero, for every choice described above.
It also contains all the chosen direction, character, and
frame coefficients. It can be chosen Galois over $\mathbb Q$.
The original field $E$, primes, Hermitian lattices, weights,
and local moduli remain the fixed data in this assertion.
\end{proposition}
\begin{proof}
For each of the finitely many positive Hermitian spaces,
choose a finite set of representatives of integral self-dual
lattices, using Theorem~\ref{al-finite-classes}. For each
pair of representatives, each of the finitely many $S,M,c$,
and each class in (JF3), choose one theta indicator. There
are finitely many such functions. In a fixed finite number
of variables there are only finitely many monomials of total
degree at most $d$. Adjoin all the corresponding algebraic
mixed derivatives, all $c_\eta$, the normal closure of $E$,
and the finite character and frame data described above.
The resulting extension of $\mathbb Q$ is finite, since
adjoining a finite list of algebraic numbers gives a finite
extension by the tower formula.

Consider any other pair $\gamma L,\gamma'L'$ with
$\gamma,\gamma'$ rational unitary transformations. The
covariance in (TE17), applied term by term to the finite
indicator, carries its variable by
\[
                z_\eta\longmapsto
                             \eta(\gamma')^{-1}z_\eta\eta(\gamma).
                                                               \tag{JF4}
\]
It carries the lift of $x$ to the lift of $\gamma'x\gamma^{-1}$;
local quotients, their indices, and the Hermitian pairing
are preserved. Thus this is equality of the actual finite
averages, not merely equality up to a complex scalar.
After dividing by the determinant character, its remaining
scalar is a product of integer powers of
$\prod_\eta\eta(\det\gamma)$ and
$\prod_\eta\eta(\det\gamma')$. The principal-ideal rule
(JF2) gives this statement with no new roots. This scalar
and every entry of the linear substitution (JF4) belong to
the normal closure of $E$.

For clarity, if $G$ is holomorphic and $A$ is a linear
substitution whose entries lie in $K$, the ordinary chain
rule gives
\[
 D_{v_1}\cdots D_{v_j}(G\circ A)(0)
                    =D_{Av_1}\cdots D_{Av_j}G(0).        \tag{JF5}
\]
Expanding each $Av_i$ in the coordinate basis expresses the
right side as a finite sum of derivatives of order $j$ of
$G$, with coefficients in $K$. The factors $c_\eta$ in the
chosen directions and their inverses are already in $K$.
Therefore every required derivative for the new pair lies
in the same field. This covers all self-dual lattices, with
no new generators as the isometries vary.

Taking the normal closure of this finite field proves the
last field assertion. Complex conjugation then restricts
to an automorphism of $K$. This does not assert that $K$
is a CM field; no such assertion is needed here.
\end{proof}

\subsection{Order of choices and descent of the column identities}

The finite generator set just constructed depends only on
the data fixed before (JF1). In particular the primes in
$\mathcal B$ are primes of $E$, not primes of the enlarged
coefficient field. Field enlargement does not replace them
by a new list or change a lattice over $E$. The finite list
of sets $S\subseteq\mathcal B$ already includes every later
allowed selection, and (JF1) already includes every required
local modulus. The derivative bounds depend on the chosen
weights and ranks, which also precede the field enlargement.
An arbitrary locally constant $K$-valued test function on
one of the finite quotients is a $K$-linear combination of
its characteristic functions. Its values therefore do not
require further generators.

In particular there is no requirement that the selected
primes remain unramified in $K$. The ordinary moment and
valuation arguments use their specified splitting in $E$
and in the imaginary quadratic field. At any place of $K$
above such a prime, extending the embedding to $\mathbb C_p$
gives the same normalized absolute value and the same unit
bounds. Adjoining coefficients may ramify that place; it
does not alter those absolute values. Hence no further
prime selection is caused by Proposition~\ref{jcf-common-field}.

All equalities of theta coefficients and of finite class
sums which were proved over $\mathbb C$ now descend to $K$:
both sides have entries in $K$, and the embedding
$K\hookrightarrow\mathbb C$ is injective. In particular
this removes the coefficient-field proviso from the Hecke
commutation identity for these fixed columns. Every
required mixed derivative belongs to the source's set of
rational indices. The assertion is uniform over the fixed
finite construction; it does not assert that one number
field contains theta jets of unbounded order or at
unboundedly many independently chosen levels.


